% Generated draft: requires mathematical and source review.
\documentclass[11pt]{article}
\usepackage[margin=25mm]{geometry}
\usepackage{fontspec}
\setmainfont{Latin Modern Roman}
\usepackage{amsmath,amssymb,mathtools,unicode-math}
\setmathfont{Latin Modern Math}
\usepackage{enumitem,needspace,xurl,hyperref}
\setlength{\parindent}{0pt}
\setlength{\parskip}{5pt}
\setlength{\emergencystretch}{3em}
\newcommand{\N}{\mathbb{N}}
\newcommand{\Z}{\mathbb{Z}}
\newcommand{\Q}{\mathbb{Q}}
\newcommand{\R}{\mathbb{R}}
\newcommand{\C}{\mathbb{C}}
\newcommand{\F}{\mathbb{F}}
\newcommand{\A}{\mathbb{A}}
\newcommand{\PP}{\mathbb{P}}
\newcommand{\E}{\mathbb{E}}
\newcommand{\eps}{\varepsilon}
\newcommand{\dd}{\,\mathrm{d}}
\begin{document}
\section*{UnsolvedMath ID: 800977}
\textbf{Source problem code:} catalogue entry
\subsection{Definitions and mathematical statement}
A closed surface is a compact surface without boundary. Under the usual smooth, strongly convex convention, a Finsler metric on a smooth surface \(M\) is a function \(F:TM\to[0,\infty)\), continuous everywhere and smooth off the zero section, such that \(F(x,v)>0\) for \(v\ne0\), \(F(x,tv)=tF(x,v)\) for \(t>0\), and the fibre Hessian of \(\tfrac12F^2\) is positive definite at every nonzero tangent vector. No reversibility assumption \(F(x,-v)=F(x,v)\) is stated in the source. Geodesics are understood through the Euler--Lagrange equations for \(\tfrac12F^2\); an unparameterized geodesic retains its curve while forgetting its parametrization.

\textbf{Question (supplied formulation).} Let \(F_1,F_2\) be Finsler metrics on a closed surface \(M\) of genus \(g\ge2\). Assume that every unparameterized \(F_1\)-geodesic is an \(F_2\)-geodesic. Must there exist a constant \(c>0\) and a smooth differential one-form \(\beta\) on \(M\), with \(\mathrm d\beta=0\), such that
\[
F_1(x,v)=cF_2(x,v)+\beta_x(v)
\qquad\text{for every }x\in M\text{ and }v\in T_xM?
\]
Here addition of a closed form means the displayed fibrewise equality, and the multiplier is a single constant on the surface.

\textbf{Unresolved qualifications.} The source does not explicitly specify orientability of the surface or the regularity and convexity convention for Finsler metrics. The definition above records the usual smooth, strongly convex reading. For possibly nonreversible metrics, the source also does not specify whether the reparametrizations implicit in ``unparameterized'' must preserve orientation. These qualifications require clarification before a fully unambiguous formulation can be accepted.
\subsection{Short English statement}
If every geodesic of one Finsler metric on a closed surface of genus at least two is also a geodesic of another after forgetting parametrization, must the first metric equal a positive constant times the second plus a closed one-form? The source leaves orientability, metric regularity, and orientation of reparametrization implicit.
\subsection{Sources}
\begin{itemize}
\item Supplied local UnsolvedMath registry record, permanent ID \(800977\), exact source problem code \texttt{catalogue entry}; title: \emph{Open Problems and Questions about Geodesics --- Question 58}. The record attributes the problem to Burns, Matveev, with proposed year 2019, and labels it \texttt{open}. This status is supplied registry data, not an independently verified literature assessment.
\item Burns, Matveev, \emph{Open Problems and Questions about Geodesics} (2019, as identified by the supplied record), Question 58. Supplied source URL: \url{https://arxiv.org/abs/1308.5417}. The background records access on 2026-10-16 and extraction by \texttt{arxiv\_latex}; these are supplied provenance claims, not independent verification of access or source contents.
\end{itemize}
\end{document}
